\chapter{Methodology}
\label{sec:methodology}

\section{The Experimental Design: Two Axes, Four Variants}
\label{sec:experimental-design}

The main experiment compares encodings of the same problem that differ in how the domain-specific heuristic is represented and interpreted. Within each family, the problem constraints are held fixed while these two properties vary. The exploratory variants separately examine auxiliary predicates and, on BSP, a compact reformulation of the balance constraint.

The \emph{approach} axis fixes where the heuristic is evaluated. In the ground-and-solve approach the heuristic is written with native \texttt{\#heuristic} directives: gringo expands them before search and clasp's Domain heuristic consumes the resulting ground objects. In the lazy approach the heuristic remains a handful of ground facts, \texttt{\_\_heuristic(...)} for the native backend and \texttt{heuristic("...")} for the Prolog one, which the propagator reads at its initialization and evaluates at every decision step against the partial assignment.

The \emph{semantics} axis fixes how a heuristic condition reads partial information. Under the clingo-like reading, \texttt{not a} requires $a$ to be assigned false; under the Alpha-like reading, it holds whenever $a$ is not assigned true. Alpha-like aggregates are computed over the atoms currently true, allowing the BSP heuristic to compare the two partial sums during search. The lazy clingo-like variant instead waits until all aggregate sources are assigned. This is the policy implemented for aggregate bindings in our evaluator; clingo itself can determine some aggregate conditions, such as monotone bounds, earlier.

\begin{figure}[H]
\centering
\includegraphics[width=0.94\textwidth]{axes}
\caption{The experimental matrix. Prefixes name the approach, suffixes the semantics.}
\label{fig:axes}
\end{figure}

Crossing the two axes yields the four variants that the rest of the thesis refers to as: \texttt{gc} and \texttt{ga} for the ground approach under the clingo and the Alpha semantics, \texttt{lc} and \texttt{la} for the lazily evaluated one, the prefix naming the approach and the suffix the semantics. To them is added the heuristic-free baseline \texttt{gc\_noheur}, which carries the same domain logic with no heuristic at all and measures what the solver does when it is left alone (Figure~\ref{fig:axes}).

To isolate specific architectural and modeling mechanisms beyond this primary matrix, the experimental design also incorporates two exploratory variants, \texttt{ga\_weak} and \texttt{la\_aux}, and a second set of BSP encodings, the \texttt{\_co} family, covering the four main variants, the heuristic-free baseline, and \texttt{ga\_weak}. Here, the \emph{BSP matrix} denotes this collection of experimental configurations for the Balanced Sum Problem. The \texttt{\_co} family reformulates the problem's balance constraint, which requires the sums of the two partitions to differ by at most one. This is a constraint defining admissible solutions, rather than a condition controlling heuristic applicability. The empirical motivation and methodological role of these exploratory variants in disentangling the contributions of default negation, body materialization, and constraint formulation are discussed in Section~\ref{sec:exploratory-methodology}.

\subsection{Simulating Alpha semantics}
\label{sec:ga-simulation}

Reproducing an Alpha-like behaviour within clingo requires addressing three distinct aspects: default negation over unassigned literals, dynamic aggregate evaluation, and the value an aggregate takes before its sources exist.

\subsubsection{Default negation over unassigned literals}

The first transformation concerns negative literals in heuristic bodies: \texttt{not c(X)} in BSP, \texttt{not not\_assigned\_sensor\_unit(S,U)} in PUP, and the guards \texttt{not fullCabinet(C)} and \texttt{not assignedThing(T)} in HRP. 

clasp applies heuristic modifications only to variables that are currently unassigned, as enforced in \texttt{DomainHeuristic::propagate}:
\begin{minted}{cpp}
if (s.value(a.var) == value_free && a.prio >= prio) { applyAction(s, a, prio); }
\end{minted}
Here, \texttt{s} is the solver and \texttt{a} is the heuristic action associated with a satisfied condition. The field \texttt{a.var} identifies its target variable; \texttt{s.value(a.var) == value\_free} checks that this variable is still unassigned. The comparison \texttt{a.prio >= prio} checks the action's priority against the priority currently stored for the same target variable and modifier. Only when both checks succeed does \texttt{applyAction} apply the heuristic modification.
It follows that \emph{a directive whose body entails the truth value of its own target atom can never fire}: once the body conditions hold, the target is already decided, and clasp skips it.

All heuristic directives in the suite that test the search state through negative literals show the same problem:
\begin{itemize}
    \item In \textbf{BSP}, \texttt{b(X)} and \texttt{c(X)} are complementary choices, so the condition \texttt{not b(X)} is satisfied only once the target \texttt{c(X)} is already derived.
    \item In \textbf{PUP}, the heuristic body repeats the negative literal of the rule that derives the target atom, namely \texttt{not not\_assigned\_sensor\_unit(S,U)}.
    \item In \textbf{HRP}, the placement directives for \texttt{cabinetTOthing(C,T)} require the guard \texttt{not assignedThing(T)}; in ground clingo this holds only when the candidate assignments for $T$ are already proven false, meaning the target is never unassigned (\texttt{value\_free}) when the directive evaluates.
\end{itemize}
Consequently, none of these directives ever fires in the \texttt{gc} variant. In HRP the same encoding also carries directives with no condition at all, which are applied at initialization and do steer the search; the conditional ones listed above stay inert.

The ground variant \texttt{ga} therefore drops these negative literals from the heuristic bodies. 
For directives where the negated literal simply tests the target atom or its complement (as in BSP and PUP), this preserves the intended applicability on unassigned targets, and clasp's \texttt{value\_free} check still prevents modifying an already assigned atom. 
In HRP, however, the dropped literals also include capacity guards such as \texttt{not fullCabinet(C)}, which describe a shared resource rather than the target atom itself. There, dropping the guard is a static approximation used by this translation: it enables the directive to fire on unassigned atoms from the start of search, at the price of never deactivating once a cabinet becomes full.

In short, deleting \texttt{not a} from a heuristic body preserves the Alpha reading under the sufficient condition that $a$ is the target atom or its complement. For other atoms, such as the HRP capacity guards, this argument does not establish equivalence; removing the guard can change which decisions the heuristic proposes. The HRP translation is therefore not established as equivalent to the Alpha-like heuristic. A treatment that preserves the guards' intended effect using clingo priorities remains future work.


\subsubsection{Dynamic aggregates by bound-guarded unrolling}

The second transformation replaces dynamic aggregate equality with bound-guarded unrolling over all candidate values. For BSP:

\begin{minted}{prolog}
sum_value(0..tot).
#heuristic b(X) : x(X), sum_value(S), S <= #sum{Y : c(Y)},
    W = ((n+1)*S)+X. [W, true]
\end{minted}

Here the directive is unrolled over all \texttt{sum\_value(S)} and guarded by a bound condition, so that during search the currently satisfied bounds encode the current aggregate value, provided the guard is monotone, i.e.\ once satisfied it stays satisfied as further atoms become true.

Without explicit priorities all unrolled directives share the default priority $0$. clasp applies newly satisfied directives by overwriting previous ones; it therefore retains the \emph{most recently activated} directive, not the maximum-weight one. To match the reference heuristic, the unrolling must align the activation order with the weight order:
\begin{itemize}
    \item \textbf{Lower-bound guards ($V \le f$):} As $f$ increases, the active set $\{V : V \le f_{\mathit{curr}}\}$ expands upwards. The last activated value is $V = f_{\mathit{curr}}$, which yields the maximum weight if $W(V)$ is strictly \emph{increasing} (e.g.\ in BSP, where $W(S) = (n{+}1)S + X$).
    \item \textbf{Upper-bound guards ($V \ge f$):} The active set $\{V : V \ge f_{\mathit{curr}}\}$ instead expands downwards, as $f$ decreases, and the alignment requires $W(V)$ to be strictly \emph{decreasing}.
\end{itemize}

\subsubsection{Empty aggregates}

\texttt{\#max} over an empty set is $0$ in Alpha but $-\infty$ (\texttt{\#inf}) in clingo, and an aggregate read on a partial assignment is empty exactly at the start of search, when the heuristic is most needed. All encodings of the suite therefore write \texttt{\#max\{0; ...\}} with an explicit neutral element, and both lazy backends implement the same rule: the native \texttt{MaxState} returns $0$ when empty, and the Prolog runtime defines \texttt{dyn\_max} with a $0$ default.

\subsection{Methodological Rationale for the Exploratory Variants}
\label{sec:exploratory-methodology}

While the four main variants evaluate the core matrix of approach versus semantics, the following are introduced as experimental controls to isolate specific factors:

\paragraph{Isolating Alpha-Negation from Dynamic Aggregates (\texttt{ga\_weak}).}
Alpha-like semantics treat default negation as satisfied over unassigned literals, and evaluate aggregates dynamically over currently true atoms. While \texttt{ga} simulates both via extensive ground unrolling, \texttt{ga\_weak} acts as an ablation study that adapts the negative guards while retaining aggregate equality bindings instead of the monotone bounds used by \texttt{ga}. This tests whether adapting default negation alone improves search on these encodings, and what further benefit comes from dynamic aggregate evaluation.

\paragraph{Fewer directives, or less grounding? (\texttt{la\_aux}).}
The \texttt{la\_aux} variant retains the lazy propagator but expresses the BSP heuristic conditions as ordinary ASP rules. For the target \texttt{b(X)}, the condition is rewritten as:
\begin{minted}{prolog}
h_b(X,S) :- x(X), not c(X), S = #sum{Y : c(Y)}.
\end{minted}
A rule defines \texttt{h\_c(X,S)}. At decision time, the lazy evaluator queries the auxiliary atoms that hold, matches their first argument to the target element $X$, and uses their second argument $S$ to compute the weight $(n+1)S+X$. Thus the conditions are instantiated during grounding and evaluated by the solver during search; grounding does not precompute their truth values for future assignments.

The encoding still contains two lazy template facts, one for each target predicate, but the auxiliary rules are grounded over candidate pairs $(X,S)$. The small template count therefore does not imply a small ground representation of the heuristic conditions. These templates must be distinguished from native \texttt{\#heuristic} directives and their ground instances. The variant explores the effect of materializing conditions while retaining lazy selection; it does not by itself establish an internal representation different from clasp's treatment of native heuristic bodies, or isolate representation cost from changes in condition applicability.

\paragraph{Disentangling Heuristic and Problem Constraint Bottlenecks (the \texttt{\_co} family).}
In the baseline BSP encoding, the balance requirement is enforced by computing the sums of both partitions independently and verifying that their difference does not exceed one via an integrity constraint (\texttt{:- sumB(SB), sumC(SC), (SB - SC)**2 > 1}). Because the possible partition sums scale quadratically with $n$, joining all pairs of candidate sums during grounding triggers a severe combinatorial explosion (materializing over $52.8$\,million propositional rules at $n=120$). Grounding the base constraints then dominates total runtime across the variants, making the cost of the heuristic layer harder to distinguish.

The compact (\texttt{\_co}) transformation resolves this bottleneck by exploiting the invariant that the total sum of all elements is known a priori ($tot = \sum_{i=1}^n i$). Instead of generating auxiliary predicates for both sums and checking them quadratically, the balance condition is restated as a single aggregate constraint with constant interval bounds, requiring the sum of partition $b$ to fall within $[\lfloor tot / 2 \rfloor, \lceil tot / 2 \rceil]$ (Section~\ref{sec:bsp}). This removes the cross-product join over candidate sums entirely, collapsing the base encoding by five orders of magnitude (from millions of rules to just a few hundred). As an experimental control, the \texttt{\_co} family eliminates the base grounding overhead, cleanly isolating the representational cost of the heuristics and clarifying whether dynamic guidance remains beneficial when domain propagation is already efficient.

\subsection{Weights, Priorities, and Flattening}
\label{sec:flattening}

Section~\ref{sec:bg-alpha} recorded that the two systems read the annotation $[W@P]$ of a heuristic differently. In Alpha, the pair is an \emph{absolute} rank over all heuristic instances, read lexicographically, with $P$ compared first: the priority $P$ is a global level and decides first, so no instance of level $P$ is considered while an instance of a level $P' > P$ is still applicable, and the weight $W$ only orders the instances that share the same level. In clingo, the evaluation is different: $P$ is local, arbitrating only between heuristics relative to the \emph{same} target atom, while the order in which distinct atoms are decided comes from the weight alone. An encoding written for Alpha therefore cannot be handed to clingo verbatim.

The intended order can nevertheless be recovered, by \emph{flattening} the level into the weight:
\[
W_{\mathit{flat}} \;=\; W + P \cdot M,
\qquad M > \max W - \min W,
\]
where the intra-level weights $W$ range over $[\min W, \max W]$ and the level multiplier $M$ is larger than that range. Under this condition $W + P\cdot M > W' + P'\cdot M$ holds whenever $P > P'$, regardless of the two individual weights, and reduces to $W > W'$ when $P = P'$: the total order induced on the atoms is then the one Alpha's absolute levels induce. When all weights are non-negative, which is the case in every encoding of this suite, $M > \max W$ suffices.

Each benchmark fixes its own $M$. In the PUP encodings $M$ is the constant offset $100{,}000$ added to the zone weights, against a maximum sensor weight of $26{,}220$ on the benchmark instances; in HRP $M$ is derived from the instance as \texttt{levelMul(LM) :- LM = MC + MR + 10} over the maximum cabinet and room identifiers, which gives $16$ on the smallest instance and $54$ on the largest. BSP does not need a multiplier, as discussed in Section~\ref{sec:bsp}.

\subsection{The Bounded Weight Field}
\label{sec:weight-problem}
clasp stores the weight of a \texttt{\#heuristic} modification in a signed 16-bit integer field and clamps any value exceeding $32{,}767$. Consequently, two ground directives whose weights both exceed $32{,}767$ become indistinguishable to clasp's Domain heuristic. A flattened weight must therefore satisfy $W + P\cdot M \le 32{,}767$ for the ground variants to rank as intended. 
HRP stays far below this bound; PUP exceeds it on the zone directives, with the consequences discussed in Section~\ref{sec:pup}; and BSP, whose maximum weight grows as $\Theta(n^{3})$, crosses it within the benchmarked range. What matters there is not the largest weight the grounder emits, but the largest one the search actually reaches: directives whose aggregate value is never attained may clamp without any observable effect. On the tested BSP sizes, the ground simulation matches the lazy search up to $n=50$, deteriorates at $n=60$, and times out from $n=70$, as reported in Section~\ref{sec:ga-weak}. 

In the compact \texttt{\_co} encoding (Section~\ref{sec:co-family}), weights saturate in exactly the same way, yet runtime is unaffected because propagation resolves choices directly without needing heuristic guidance.
The lazy backends avoid this limitation altogether. Rather than assigning static weights to clasp, the propagator ranks candidate targets dynamically at runtime and directly returns the chosen decision literal.

The suite folds each benchmark's global levels into the weights in all four variants and both backends, so cross-atom ordering does not depend on backend-specific priority semantics. Accordingly, the native lazy language (Section~\ref{sec:dsl}) omits a separate priority field. The Prolog representation retains one for interface compatibility: BSP and PUP use \texttt{P=0}, whereas HRP assigns \texttt{P=1} to its nine ranking rules and \texttt{P=0} to the final false layer, as detailed in Section~\ref{sec:dsl-limits}.

\section{System Architecture}
\label{sec:architecture}

The project maintains two modified copies of clingo. They share the same upstream sources and differ only in their extensions: the build in \path{clingo-native} carries the C++ backend, while the build in \path{clingo-prolog} carries the Prolog backend, configured with \texttt{clingo\_USE\_SWIPL=ON} so that the SWI-Prolog engine is linked in-process, making its memory impact easier to account for. In both builds, the extension lives entirely under \path{libclingo}, exposes the same class name, and is registered at the same point of \path{clingo_app.cc}:

\begin{minted}{cpp}
HeuristicPropagator my_propagator;
ctl.register_propagator(my_propagator, true);
\end{minted}

The propagator is registered on \emph{every} run of the modified binary. On a program that contains no lazy heuristic facts, the propagator itself is inactive: \texttt{init} finds nothing to watch and \texttt{decide} always returns $0$, so every decision is left to clasp's own branching heuristic. The same binary therefore also runs the ground variants (with clingo's standard \texttt{\#heuristic}s) and the heuristic-free baseline without requiring a separate build. 

In the registration call \texttt{ctl.register\_propagator(my\_propagator, true)} shown above, the second argument, \texttt{true}, instructs clasp to call the propagator sequentially, avoiding race conditions on the incremental state described in Section~\ref{sec:runtime}. No other parts of clingo are modified. Figure~\ref{fig:architecture} shows the architecture of the two binary build variants.

\begin{figure}[H]
\centering
\includegraphics[width=\textwidth]{architecture}
\caption{System architecture of the two binary variants, showing their respective heuristic propagators and evaluation backends attached to clasp.}
\label{fig:architecture}
\end{figure}

In both builds, the class implements \texttt{clingo::Heuristic}, the propagator interface extended with the \texttt{decide} callback (Section~\ref{sec:bg-propagator}). Both backends follow the same lifecycle, built on \texttt{init}, \texttt{propagate}, \texttt{undo} and \texttt{decide}:
\begin{itemize}[leftmargin=1.5em, itemsep=1pt]
    \item During \texttt{init}, they scan symbolic atoms for trigger facts (\texttt{\_\_heuristic(...)} in the C++ build and \texttt{heuristic("...")} in the Prolog build), construct their runtime representation, and register watches on solver literals.
    \item During forward search, \texttt{propagate} updates the internal state of the aggregates appearing in the heuristic bodies, together with the candidate states, as watched literals become true.
    \item During backtracking, \texttt{undo} restores those aggregate and candidate states as literals are retracted from the trail.
    \item At each decision point, \texttt{decide} either returns a signed literal for the highest-ranked target or returns $0$, delegating the decision to clasp's fallback heuristic.
\end{itemize}
The two backends support different heuristic languages. The native backend maintains an incremental state for \texttt{\#sum}, \texttt{\#count}, \texttt{\#min} and \texttt{\#max} over one predicate and one argument position, and can only filter the aggregated set by matching its arguments with those of the target atom; the Prolog backend evaluates an arbitrary \texttt{aggregate\_all/3} goal and is unrestricted in this respect. Section~\ref{sec:dsl} states the language subset each backend accepts, and Section~\ref{sec:dsl-limits} what falls outside it.

To activate this mechanism, runs are launched with the flag \texttt{--heuristic=Domain}. A detailed discussion of the architectural differences between the two backends is provided in Chapter~\ref{sec:implementation}.

The propagator never adds a constraint, never derives a literal, and never refuses one. It changes only the order and preferred polarity of decision literals, so the answer sets are untouched; a heuristic can make a run slower or faster but cannot change the meaning of the encoding (detailed in Section~\ref{sec:decision}).

\section{Translation of the Heuristic Directives}
\label{sec:dsl}

Neither backend reads native \texttt{\#heuristic} directives directly. Instead, a custom declarative language replaces each directive with a single ground fact that \emph{describes} the heuristic rule.

This substitution is driven by two main reasons:
\begin{itemize}[leftmargin=1.5em, itemsep=2pt]
    \item gringo treats a \texttt{\#heuristic} directive as an ordinary ASP rule and immediately grounds it. By the time the propagator initializes in \texttt{init}, the combinatorial explosion has already occurred. What the propagator needs is the high-level \emph{template}, which gringo never grounds and the clingo C++ API does not expose; writing it as an ordinary ground fact is what makes it reachable from \texttt{init}.
    \item Aggregates must survive grounding unevaluated so they can be monitored dynamically during search. Standard ASP syntax provides no such mechanism: an aggregate inside a rule body is either evaluated statically by the grounder or unrolled over its value domain.
\end{itemize}

The lazy fact is therefore the original directive rewritten in a form that passes through the grounder untouched. Its structure mirrors the original construct-by-construct. 
The remainder of this section details this representation. We introduce each construct by rewriting the BSP heuristic (Section~\ref{sec:motivation}) step by step. Finally, Section~\ref{sec:dsl-limits} outlines the expressive limitations of both backends.

\subsection{Lazy heuristic syntax for the C++ backend}

\begin{minted}{prolog}
#heuristic b(X) : x(X), not c(X), S = #sum{Y : c(Y)}, W = ((n+1)*S)+X. [W, true]
\end{minted}

\paragraph{Target.} The target atom is \texttt{b(X)}, where $X$ ranges over the ground instances. The lazy fact specifies only the predicate name, \texttt{\_\_target(b)}. During \texttt{init}, the propagator queries clingo's atom table for all ground instances of \texttt{b(...)} already materialized by gringo.

\paragraph{Positive body.} The condition \texttt{x(X)} shares its argument with the target: the candidate for \texttt{b(7)} is alive only while \texttt{x(7)} is true. A condition of this shape is written as the bare predicate name, \texttt{x}, and the propagator matches its arguments against the target tuple position by position.

If a body predicate has a different arity or argument order, or supplies a value used in the weight, the correspondence must be stated explicitly:

\begin{minted}{prolog}
__body(pred, __match(src, tgt), __bind_arg(var, idx))
\end{minted}

\texttt{\_\_match(src, tgt)} requires the argument at the \texttt{src} position of the body atom to equal the argument at the \texttt{tgt} position of the target atom, where both \texttt{src} and \texttt{tgt} are integers. All indices are 0-based, and at least one \texttt{\_\_match} is mandatory. \texttt{\_\_bind\_arg(var, idx)} extracts argument \texttt{idx} of the body atom into a variable.

Given the ground facts:
\begin{minted}{prolog}
score(1,10). score(2,5). score(3,7).
\end{minted}
and the target atoms:
\begin{minted}{prolog}
{ choose(1..3) }.
\end{minted}
the directive:
\begin{minted}{prolog}
__heuristic(__target(choose), __body(score, __match(0,0), __bind_arg(w,1)), __weight(w), __modifier(true)).
\end{minted}

pairs each \texttt{choose(X)} with the \texttt{score(X,W)} that agrees with it on the first argument, and binds \texttt{w} to the second argument of that atom. The weights are then $10$, $5$ and $7$, so the propagator decides \texttt{choose(1)} first.

\paragraph{Negative body.} The condition \texttt{not c(X)} likewise matches its arguments against the target tuple. Negative body literals are specified by prefixing the predicate name with \texttt{\_\_n\_}.

\paragraph{Semantics.} The reading of that negation is fixed by the \texttt{\_\_semantics(S)} selector, which decides whether \texttt{\_\_n\_c} means \enquote{\texttt{c(X)} is assigned false} or \enquote{\texttt{c(X)} is not assigned true}. The same selector fixes the reading of aggregates, described below. It takes \texttt{alpha} or \texttt{clingo}, and defaults to \texttt{clingo}, selecting the clingo-like reading of body conditions described above.

\paragraph{Aggregates.} Aggregates are evaluated lazily during search. For instance, \texttt{\_\_bind(s, \_\_sum(c))} binds variable \texttt{s} to the sum over the predicate \texttt{c} (by default aggregating its first argument). The operators \texttt{\_\_count}, \texttt{\_\_min}, and \texttt{\_\_max} follow the same structure.

An aggregate can be restricted using \texttt{\_\_filter(src\_idx, tgt\_idx, offset)}, which narrows it \emph{per target}, keeping only source atoms whose argument at position \texttt{src\_idx} equals argument \texttt{tgt\_idx} of the target plus an optional \texttt{offset}. For instance, in PUP, retrieving the maximum number of sensors on the \emph{previous} unit ($U-1$) for a target \texttt{assigned\_zone\_unit(Z,U)} is expressed as:
\begin{minted}{prolog}
__bind(m1, __max(num_sensors_on_unit, 0, __filter(1, 1, -1)))
\end{minted}
where \texttt{0} selects the aggregated argument $N$ in \texttt{num\_sensors\_on\_unit(N, U')}, and the filter matches $U' = U - 1$.

\paragraph{Weight.} Weight is expressed as a \emph{term tree} such as:
\begin{minted}{prolog}
__weight(__add(__mul(s, B), self))
\end{minted}
which represents $(s \cdot B + \mathit{self})$, where \texttt{self} is a shorthand for the target atom's argument at index 0 without requiring a variable binding.

The binary constructors \texttt{\_\_add}, \texttt{\_\_sub}, and \texttt{\_\_mul} form expression trees whose leaves belong to three categories: integer constants, variables of the heuristic language, and target arguments (accessed through \texttt{self} or \texttt{\_\_bind\_target\_arg(var, idx)} for any other).

At \texttt{init}, the propagator compiles all weight term trees into ASTs. During search, each active heuristic object evaluates its weight AST dynamically against the current trail state, avoiding the memory overhead of materializing them in the ground program. Only candidates with a strictly positive resolved weight are taken into account; the others are ignored, leaving their decision to the solver's default heuristic.

\paragraph{Modifier.} The \texttt{\_\_modifier(M)} selector specifies which of clingo's heuristic modifiers to use: \texttt{level} sets the ranking value alone, \texttt{sign} sets the preferred polarity alone, while \texttt{true} (the default) and \texttt{false} set both. 

The solver's \texttt{decide} interface expects a \emph{signed} literal, so a ranking value alone does not determine a decision. Each target keeps one slot per channel, ranking value and preferred polarity, and when several directives write to the same slot the one with the highest weight wins. This is what replaces clingo's local priority. Across the benchmarks all directives use \texttt{true}, except for the final layer of HRP, which uses \texttt{false} to prune the remaining choice atoms.

\medskip
Assembling the components gives the fact in full:

\begin{minted}{prolog}
__heuristic(__target(b), 
            x, 
            __n_c, 
            __bind(s, __sum(c)), 
            __weight(__add(__mul(s, B), self)), 
            __semantics(alpha), 
            __modifier(true)) 
                :- B = n + 1.
\end{minted}


\subsection{Lazy heuristic syntax for the Prolog backend}
\label{sec:prolog-notation}

Unlike the native backend, which supports a fixed heuristic DSL, the Prolog backend accepts standard goals available to the embedded engine. A heuristic is declared as a string inside a \texttt{heuristic("...")} fact and parsed directly as a Prolog clause.

\begin{minted}{prolog}
heuristic(b(X), W, 0, true) :-
    holds(x(X)),
    target_available(b(X)),
    alpha_not(c(X)),
    aggregate_all(sum(Y), holds(c(Y)), S),
    holds(heuristic_base(Base)),
    W is Base * S + X.
\end{minted}

The \textbf{head} specifies the target atom ($A$), weight ($W$), priority ($P$) and modifier ($M$). The \textbf{body} gives the conditions and computations used to derive them. Priorities follow the convention described in Section~\ref{sec:weight-problem}.

To interact with the solver's search state, the backend provides a few dedicated interface predicates: \texttt{holds/1} succeeds on atoms currently true on the trail and on static ground facts; \texttt{alpha\_not/1} and \texttt{clingo\_not/1} implement the two readings of default negation described in Section~\ref{sec:bg-alpha}; and \texttt{target\_available/1} discards candidates that are already assigned.

Everything else uses standard Prolog: arithmetic relies on \texttt{is/2}, and set aggregations use \texttt{aggregate\_all/3}. The backend also provides the wrappers \texttt{dyn\_sum/3}, \texttt{dyn\_count/2}, \texttt{dyn\_min/3}, and \texttt{dyn\_max/3}. When no matching atoms are assigned on the trail, \texttt{dyn\_min/3} and \texttt{dyn\_max/3} return $0$ instead of failing. This lets the heuristic compute a weight even at the start of search.

Table~\ref{tab:translation} summarizes the correspondence between the three notations.


\begin{table}[H]
\centering
\footnotesize
\setlength{\tabcolsep}{4pt}
\renewcommand{\arraystretch}{1.15}
\begin{tabular}{@{}>{\raggedright\arraybackslash}p{2.1cm}>{\raggedright\arraybackslash}p{3.3cm}>{\raggedright\arraybackslash}p{3.9cm}>{\raggedright\arraybackslash}p{3.9cm}@{}}
\toprule
\textbf{construct} & \textbf{clingo directive} & \textbf{native lazy fact} & \textbf{Prolog rule string}\\
\midrule
\multicolumn{4}{@{}l}{\emph{structure}}\\
\addlinespace[2pt]
target & \texttt{b(X)} & \texttt{\_\_target(b)} & \texttt{heuristic(b(X), W, P, M)}\\
\addlinespace[3pt]
positive body, target tuple & \texttt{x(X)} & \texttt{x} & \texttt{holds(x(X))}\\
\addlinespace[3pt]
positive body, join & \texttt{p(X,V)} & \texttt{\_\_body(p,} \newline \texttt{\ \ \_\_match(0,0),} \newline \texttt{\ \ \_\_bind\_arg(v,1))} & \texttt{holds(p(X,V))}\\
\addlinespace[3pt]
negation & \texttt{not c(X)} & \texttt{\_\_n\_c} & \texttt{alpha\_not(c(X))} \newline \texttt{clingo\_not(c(X))}\\
\addlinespace[3pt]
aggregate & \texttt{S = \#sum\{Y : c(Y)\}} & \texttt{\_\_bind(s,} \newline \texttt{\ \ \_\_sum(c))} \newline also \texttt{\_\_count}, \texttt{\_\_min}, \texttt{\_\_max} & \texttt{aggregate\_all(} \newline \texttt{\ \ sum(Y), holds(c(Y)), S)}\\
\addlinespace[3pt]
filtered aggregate & \texttt{M = \#max\{N :} \newline \texttt{\ \ q(N,U-1)\}} & \texttt{\_\_bind(m,} \newline \texttt{\ \ \_\_max(q, 0,} \newline \texttt{\ \ \ \ \_\_filter(1,1,-1)))} & \texttt{U1 is U-1,} \newline \texttt{dyn\_max(} \newline \texttt{\ \ holds(q(N,U1)),} \newline \texttt{\ \ N, M)}\\
\addlinespace[3pt]
target arguments & \texttt{X}, \texttt{Y} in \texttt{b(X,Y)} & \texttt{self}, \newline \texttt{\_\_bind\_target\_arg(} \newline \texttt{\ \ y, 1)} & ordinary head variables\\
\addlinespace[4pt]
\midrule
\multicolumn{4}{@{}l}{\emph{annotations}}\\
\addlinespace[2pt]
weight & \texttt{W = ((n+1)*S)+X} & \texttt{\_\_weight(\_\_add(} \newline \texttt{\ \ \_\_mul(s,B), self))} & \texttt{W is B*S+X}\\
\addlinespace[3pt]
modifier & \texttt{[W, true]} & \texttt{\_\_modifier(true)}, \newline default \texttt{true} & fourth head argument\\
\addlinespace[3pt]
semantics & (fixed by the system) & \texttt{\_\_semantics(alpha)} \newline \texttt{\_\_semantics(clingo)} & choice of \texttt{alpha\_not} / \texttt{clingo\_not}\\
\addlinespace[3pt]
target still free & (implicit in clasp) & implicit: an assigned target is skipped & \texttt{target\_available(} \newline \texttt{\ \ b(X))}\\
\bottomrule
\end{tabular}
\caption{The same heuristic in the three notations.}
\label{tab:translation}
\end{table}

\subsection{What the Two Languages Cannot Express}
\label{sec:dsl-limits}

While Table~\ref{tab:translation} summarizes the syntax supported by each backend, the two engines do not cover the same subset of clingo's heuristic language, nor does either support standard ASP heuristics in full. Instead, they reflect a clear design trade-off between evaluation speed and relational expressiveness.

\paragraph{The C++ language subset.}
The C++ backend is designed for high performance through incremental state tracking, but this restricts its language to a minimal template. Directives must target entire predicates rather than individual ground atoms, and arguments are strictly limited to numeric tuples. 

In rule bodies, literals can only link directly to the target tuple: arbitrary joins between body atoms (such as \texttt{p(X,V), q(V,W)}) and relational comparisons (\texttt{N < M}) are unsupported, requiring precomputation via auxiliary ASP rules. 
Similarly, aggregates operate strictly over a single predicate and argument, without support for internal joins or conditions. 

Finally, weight arithmetic is restricted to basic operations ($+$, $-$, $\cdot$) yielding strictly positive values, and modifiers altering internal VSIDS scores (\texttt{init}, \texttt{factor}) are rejected. Because the backend also omits an explicit priority field, multi-level priority hierarchies must be flattened into weights (Section~\ref{sec:flattening}).


\paragraph{The Prolog language subset.}
In contrast, the Prolog backend prioritizes expressive flexibility. Because rule bodies are standard Prolog goals, arbitrary joins, comparisons, nested aggregations (via \texttt{aggregate\_all/3}), and non-numeric terms are supported natively. Its head also carries an explicit priority argument, so it \emph{can} express Alpha's two-dimensional rank without weight flattening; the suite uses flattened weights throughout, and the priority field only in HRP, where it keeps the nine ranking rules one level above the final \texttt{false} layer. The primary limitation is that Prolog goals may generate terms that do not exist in the ground program, requiring the backend to filter out invalid candidates against clasp's atom tables (Section~\ref{sec:decision-prolog}). Unrestricted goals are re-evaluated from scratch at every decision point, introducing runtime overhead that is evaluated in Section~\ref{sec:per-decision-cost}.
